% ============================================================
%  弦紧致化的几何工具（TASI 2017 讲座）—— 中文译注本
%  TASI Lectures on Geometric Tools for String Compactifications
%  Lara B. Anderson, Mohsen Karkheiran
%  arXiv:1804.08792v1 [hep-th], 24 Apr 2018; PoS (TASI 2017)
%
%  编译：xelatex main.tex  （两遍；见 build.ps1）
%  中文字体：ctex 的 fandol 字库（随 TeX Live 分发，无需系统字体）
%  说明：导言区在 preamble.tex；正文分片在 sections/ 下。
% ============================================================
\documentclass[11pt,a4paper,UTF8,fontset=fandol]{ctexart}

% ============================================================
%  preamble.tex —— 全项目共用导言区（main.tex 与分片测试均 \input 本文件）
%  在 \documentclass[11pt,a4paper,UTF8,fontset=fandol]{ctexart} 之后使用
% ============================================================

% ---------- 页面与字体 ----------
\usepackage[a4paper,top=2.6cm,bottom=2.6cm,left=2.6cm,right=2.6cm]{geometry}
\usepackage{amsmath,amssymb,amsfonts,amsthm}
\usepackage{graphicx}
\usepackage{float}
\usepackage{caption}
\captionsetup{font=small,labelsep=quad}
\usepackage{booktabs}
\usepackage{longtable}
\usepackage{array}
\usepackage{multirow}
\usepackage{enumitem}
\usepackage{xcolor}
\usepackage{tcolorbox}
\tcbuselibrary{breakable}
\usepackage{fancyhdr}
\usepackage[colorlinks=true,linkcolor=blue!40!black,citecolor=blue!40!black,urlcolor=blue!40!black,
            bookmarksnumbered=true,bookmarksopen=true]{hyperref}
\hypersetup{
  pdftitle={弦紧致化的几何工具（TASI 2017 讲座）——中文译注本},
  pdfauthor={Lara B. Anderson, Mohsen Karkheiran（原著）；中文译注本},
  pdfsubject={弦紧致化、Calabi--Yau 流形、全纯向量丛、模空间},
  pdfkeywords={string compactification, Calabi-Yau, heterotic string, vector bundles, moduli, TASI 2017}
}

% ---------- 版式 ----------
\setlength{\parskip}{0.35ex}
\setlength{\emergencystretch}{3em}
\setlist{itemsep=0.2ex,topsep=0.5ex}
\pagestyle{fancy}
\fancyhf{}
\fancyhead[L]{\small 弦紧致化的几何工具（TASI 2017 讲座）}
\fancyhead[R]{\small\thepage}
\renewcommand{\headrulewidth}{0.4pt}

% ---------- 本节重点（Key Takeaways）----------
\newtcolorbox{keytakeaway}{
  title={本节重点（Key Takeaways）},
  fonttitle=\small\bfseries,
  coltitle=black,
  colbacktitle=gray!18,
  colback=gray!6,
  colframe=gray!45,
  boxrule=0.5pt,
  arc=1.5pt,
  left=6pt,right=6pt,top=4pt,bottom=4pt,
  fontupper=\small,
  breakable
}
\newcommand{\takeawaylabel}{\textbf{本节重点}\quad}

% ---------- 图表路径 ----------
\graphicspath{{figures/}}

% ---------- 常用符号（不改变原文记号，仅便于录入）----------
\newcommand{\Oh}{\mathcal{O}}
\newcommand{\bbZ}{\mathbb{Z}}
\newcommand{\bbC}{\mathbb{C}}
\newcommand{\bbP}{\mathbb{P}}
\newcommand{\bbR}{\mathbb{R}}
\newcommand{\bpa}{\bar\partial}


% ---------- 标题信息 ----------
\title{\bfseries 弦紧致化的几何工具\\[0.35em]
       \large —— TASI 2017 讲座 ——\\[0.8em]
       {\normalsize\normalfont$|$ 中文译注本 $|$}}
\author{Lara B. Anderson \and Mohsen Karkheiran\\[0.3em]
        {\small Virginia Tech, Department of Physics}\\[0.3em]
        {\small E-mail: [邮箱隐去], [邮箱隐去]}}
\date{}

\begin{document}

% ================= 标题页 =================
\begin{titlepage}
\centering
\vspace*{1.2cm}
{\LARGE\bfseries 弦紧致化的几何工具}\\[0.5cm]
{\large —— TASI 2017 讲座 ——}\\[1.2cm]
{\Large\itshape TASI Lectures on\\[0.2em] Geometric Tools for String Compactifications}\\[1.5cm]
{\large Lara B. Anderson \quad\&\quad Mohsen Karkheiran}\\[0.4cm]
{\normalsize Virginia Tech, Department of Physics}\\[0.4cm]
{\small E-mail: [邮箱隐去], [邮箱隐去]}\\[2.0cm]
{\normalsize 原载：arXiv:1804.08792v1 [hep-th], 24 Apr 2018}\\[0.2cm]
{\normalsize 会议：Theoretical Advanced Study Institute Summer School 2017}\\[0.1cm]
{\normalsize ``Physics at the Fundamental Frontier'', 4 June -- 1 July 2017, Boulder, Colorado}\\[1.8cm]
{\normalsize 中文译注本}\\[0.3cm]
{\small 译者说明见下页；原文版权归作者所有（CC BY-NC-ND 4.0），本译文仅供学习研究。}
\vfill
\end{titlepage}

% ================= 译者说明 =================
\thispagestyle{plain}
\section*{译者说明}
\addcontentsline{toc}{section}{译者说明}

\begin{enumerate}[label=\arabic*.,leftmargin=2em]
  \item \textbf{底本}：本译本的底本为 Lara B. Anderson 与 Mohsen Karkheiran 所著
        \textit{TASI Lectures on Geometric Tools for String Compactifications}
        （arXiv:1804.08792v1 [hep-th]，PoS 会议文集，原文共 61 页）。
        需要特别说明：本文作者是 Anderson 与 Karkheiran，并非 Edward Witten；
        正文中仅在 Gukov--Vafa--Witten 超势、Green--Schwarz--Witten 教材等确有关联处出现 Witten 的名字。

  \item \textbf{翻译原则}：逐段直译原意，不删段、不并段、不做"总结式改写"；
        术语以中文学界约定俗成者为准（如 gauge theory 译"规范场论"、holonomy 译"和乐群"、
        anomaly 译"反常"）。首次出现的术语在括号内给出英文原词。

  \item \textbf{公式}：所有带编号公式的编号与原文逐一对应（如式 (2.17)、(A.13)），
        用 \verb|\tag| 锁定，不随排版漂移；上下标、指标、希腊字母逐页对照原始 PDF 页面图像录入。
        原文笔误照录并出注，未作静默改动。

  \item \textbf{批注}：源码中以 \verb|% [Tutor Note: ...]| 形式嵌入的审校批注共一百一十余条
        （其中易错预警与含糊点拨类 74 条，"原文疑误" 28 条，"存疑" 8 条，"背景补充" 1 条；
        多行批注按末行统计，实际条数略多），覆盖易错预警（符号约定、指标升降、张量缩并）、
        含糊点拨（作者以 ``it is obvious'' 省略的证明及其依据的定理）与几何/物理直觉三类。
        这些批注为 LaTeX 注释，\emph{不进入排版结果}，在源码中可直接检索 \verb|[Tutor Note|；
        全部批注另汇总于 \verb|notes/annotations-index.md|。

  \item \textbf{图表}：原文全部 10 幅插图已逐幅从 PDF 中剥离（\verb|figures/fig01.png|---\verb|fig10.png|），
        标题译为中文并标注原文位置；表 1、表 2 按原文内容用 \verb|booktabs| 重建。

  \item \textbf{编译}：使用 XeLaTeX 编译（\verb|xelatex main.tex| 两遍），
        中文使用 ctex 的 fandol 字库，无需安装系统字体；若本机有宋体/黑体，
        可将文档类选项改为 \verb|fontset=windows|。
        pdfLaTeX 无法处理中文，请勿使用。

  \item \textbf{性质}：本译文由 AI 辅助生成并经逐页核对与编译验证，属于学习用参考译本；
        凡译文与原文有出入之处，一切以原文为准。原文采用 CC BY-NC-ND 4.0 许可。
\end{enumerate}

\vspace{1.2em}
\noindent\textbf{阅读提示}：每节开头的“本节重点”框为该节内容提要；
源码中每页起始处有以“原文 PDF 第 N 页”标记的页锚注释，可与原文逐页对照
（本 PDF 的印刷页码 $=$ PDF 页码 $-$ 1）。

\clearpage
\tableofcontents
\clearpage

% ================= 正文 =================
\input{sections/s01-intro}
\input{sections/s02-heterotic}
\input{sections/s03a-cy}
\input{sections/s03b}
\input{sections/s03c}
\input{sections/s04-bundles}
\input{sections/s05-conclusions}
\input{sections/s06-refs}

\end{document}
